\chapter{Humanismo y relectura de la Antigüedad}
El Renacimiento se caracteriza por la recuperación crítica de los textos clásicos y la reevaluación de las virtudes cívicas asociadas a la república romana y a las ciudades-estado italianas. Kemp subraya que esta relectura funcionó como un proyecto educativo que pretendía moldear tanto la vida pública como la vida privada de las elites urbanas.\cite{Kemp2008} La antología editada por Elmer y Webb muestra, además, cómo los humanistas compilaron cartas, discursos y tratados para difundir un ideal de ciudadanía letrada.\cite{Elmer2000}

El humanismo no se limitó a Italia. Popov documenta su traducción institucional en la Europa oriental, donde los principios de administración racional se injertaron sobre tradiciones locales, demostrando que el movimiento tuvo adaptaciones diversas según las condiciones políticas y lingüísticas.\cite{Popov2015}

\begin{figure}[H]
    \centering
    % Reemplazar por la imagen recortada correspondiente
    \includegraphics[width=0.75\textwidth]{FIGURES/humanistas_manoscritos_placeholder}
    \caption{Copistas trabajando sobre manuscritos clásicos ("[Streams of History Book 6] Ellwood W. Kemp - Streams of History_ The Renaissance and Reformation (2008, Yesterday's Classics) - libgen.li.epub", p. XX).}
    \label{fig:humanismo-manuscritos}
\end{figure}

La transmisión de los textos clásicos exigió talleres de copia, correctores y una economía del pergamino y del papel que antecedió a la imprenta. Ese esfuerzo técnico reforzó la idea de que el conocimiento era un capital colectivo que debía preservarse y ampliarse mediante bibliotecas y círculos de estudio.
